\section{Introduction}
\label{sec:intro}

An agent that calls tools will get some of those calls wrong. What happens next
is decided not by the model but by the harness: it appends the failed call to
the transcript, appends whatever the runtime printed, and asks the model to
continue. ReAct \citep{yao2023react} introduced this pattern and essentially
every agent framework since has kept it, including the ones built for
small open-weight models. The reasoning is hard to argue with. The error message
is real information, it was produced by ground truth rather than by a critic,
and a model that reads it should not try the same thing again.

We measured whether that is what happens. It is not.

Take a decision point in an agent trajectory: a context $\ctxpre$, an action
$\afail$ that will fail there, and the error the runtime returns. Compare the
model's probability of writing $\afail$ before the attempt with its probability
after the failure has been recorded. For every one of the
\NumModelsToolshed\ instruction-tuned models we tested, spanning
\SmallestModel--\LargestModel\ parameters across \NumFamilies\ families, that
probability goes \emph{up}, by \MeanG\ nats, or about \GPerTokenTool\ nats per
token of the action, which is a factor of \OddsPerToken\ in the odds of every
token. The effect is not an average over a few dramatic cases: it holds on
\FracNegTool\ or more of individual items.

Restated on a scale that is easier to hold onto: over a fixed set of four
candidate actions, the probability of re-emitting the failed call rises from
\PRepeatPreMean\ before the attempt to \PRepeatPostMean\ after it. Greedy
decoding, which is what the agent would actually have written next, reproduces
the failed call token for token on \GreedyPostMean\ of items after the failure,
against \GreedyPreMean\ before it.

Practitioners know the symptom. Agents get stuck repeating a call that has
already failed, and frameworks ship loop detectors and step caps to contain it.
The symptom is usually attributed to the model: it is small, it does not really
understand the error, a stronger model would recover. Our second result is that
this diagnosis is wrong, and it matters because it points at the wrong fix.

Appending $(\afail, o^{\times})$ does two things at once, and they can be
separated. It puts the token sequence of the failed call into the context, where
the model's copying machinery can reach it
\citep{olsson2022induction,xu2022breaktheloop}; and it asserts that the call
failed. We can measure each part on its own by scoring the same action under a
context whose observation is \emph{valence-free}, meaning the runtime
acknowledges the call and says nothing about how it went. The surface-form term
accounts for \CopyShare\ of the effect. The semantic term, the part that would
have to be large and negative for ``the model does not understand the error'' to
be the explanation, is small, and its sign is not even consistent between our
two environments. On tool calls it is slightly positive, since explaining what
went wrong requires naming the tool and the argument again; on program repair it
is slightly negative, so a traceback does push the model away from the code that
produced it. We stress the second case, because it is the one that fixes the
interpretation: these models are not simply failing to read error messages.
Whatever reading happens is being outvoted, by a median of \MedianCopySemRatio\
to one across model and environment, and by \MaxCopySemRatio\ to one at the
extreme.

This changes what a fix should look like. If the model misunderstood, the remedy
is a better model or a clearer message. If the harm is carried by the presence of
a string, then message engineering cannot reach it, and neither can an
instruction: telling a model not to repeat a call requires it to represent a
negation about a string that is sitting right there in its context, which is the
condition under which negative instructions are known to work least well
\citep{castricato2024pinkelephants,negation2025pinkelephant}. What should work
instead is structural.

``Structural'' turns out to be narrower than it first appears, and the
experiment that shows this is the one we did not expect. The obvious structural
fix is to delete the failed attempt and retry from a clean context, which is what
prior work on context contamination recommends \citep{yang2026contamination}.
That is the \emph{worst} harness we tested, multiplying the exact repeat rate by
\DropRepeatRatio\ without changing task success. The reason is embarrassingly
simple once seen: deleting
the failure restores the exact context that produced it, and a deterministic
policy in an identical context emits an identical action. The requirement is
therefore not ``remove the failed action'' but something more specific: the
context after a failure must differ from the context before it, and the
difference must not be the failed action.

We test all three. Replacing the verbatim call with a description the runtime
generates from its own error metadata, which keeps the diagnosis and drops the
token sequence,
removes \AbstractRemoved\ of the inversion and drives the exact greedy repeat
rate to zero, at no cost in generated tokens. A decoder-level ban on
previously-failed strings acts on the same term and costs nothing at all. The
natural-language prohibition moves the measured quantity slightly the
\emph{wrong} way, which is what the negative-instruction literature would
predict for a prohibition that has to name the very string it is prohibiting
\citep{castricato2024pinkelephants,negation2025pinkelephant}, and in
free-running rollouts it does not reduce repetition either, though it does
change behaviour in another way we report rather than explain.
Section~\ref{sec:agent} gives all of this in full, including where the
interventions do not help.

We are not claiming that execution feedback is useless, or that agent
transcripts should be thrown away. The failed call carries diagnostic
information a harness genuinely needs. The claim is narrower and, we think, more
useful: the standard way of \emph{delivering} that information hands the model a
strong reason to repeat itself, and a harness can keep the diagnosis while
withholding the reason.

\paragraph{Contributions.}
\begin{enumerate}\itemsep2pt
\item \textbf{A measurement.} We define the \emph{corrective gain} of a failure
record, the change in log-probability of re-emitting the failed action, and
show it is negative across \NumModelsToolshed\ models,
\NumFamilies\ families and two environments (simulated tool calling and MBPP
program repair with real interpreter output). We report it alongside two
interpretable readouts, a normalised repeat probability and an exact greedy
repeat rate.
\item \textbf{A decomposition that identifies the cause.} Counterfactual
observations split the gain into a surface-form term and a semantic term. The
surface-form term dominates at every scale we tested, and it, not the semantic
term, is what improves with model size.
\item \textbf{Interventions that follow from the decomposition, and one that
does not.} Replacing the verbatim failed call with a runtime-generated
description of the failure, and banning previously-failed strings at the
decoder, both act on the term that matters; adding a natural-language
prohibition does not. We evaluate all three in free-running rollouts and report
their cost.
\item \textbf{A reproducible, CPU-only artefact.} Every number comes from code
in the accompanying repository, which runs end to end without a GPU. Raw
per-score records, item sets, seeds and environment metadata are included, and
the manuscript's inline numbers are generated from them rather than typed.
\end{enumerate}
